\chapter{Literature Review}

This chapter reviews the body of knowledge relevant to the design of ETL
pipelines in regulated industries and the application of DevSecOps principles.
It begins by outlining the context of ETL and governance, defines the key
concepts central to this research, and then examines the literature associated
with each research question.

\section{Background}
\noindent
Data pipelines underpin modern analytical systems, forming the connective
infrastructure between data sources and analytic platforms
\cite{elbaghazaoui_integrating_2025}. They transform raw data into governed,
traceable, and auditable outputs used in applications such as fraud detection,
market analysis, compliance reporting, and metadata processing.

Although the extract–transform–load (ETL) process appears simple in theory, it
is in practice complex, resource-intensive, and often the most time-consuming
element of data warehousing \cite{kimball_data_2004}. ETL pipelines act as the
“eyes of the system,” feeding reliable data downstream; when lineage is lost or
visibility fails, the reliability of the entire process is compromised
\cite{elbaghazaoui_data_2025-1}. Vayyala~ \cite{elbaghazaoui_data_2025-1}
emphasises that maintaining end-to-end lineage across transformations is vital
for sustaining trust in analytics-driven decisions.

In regulated industries, compliance depends as much on how data is processed as
on the outcomes produced. Frameworks such as GDPR, SOX, and APRA’s CPS standards
require secure, ethical, and well-documented data handling
\cite{elbaghazaoui_data_2025,cps230,cps510,cps234,australian_prudential_regulation_authority_prudential_2013}.
These frameworks place accountability, ownership,
and audit trails at the centre of compliance culture.

Auditability, in this context, is the ability to demonstrate compliance through
reproducibility, lineage, and tamper-resistant logs
\cite{elbaghazaoui_ensuring_2025}. Trustworthy analytics rely not only on
integrity of data but also on immutable evidence of correct process execution.
Governance provides the organisational scaffolding for these assurances—linking
policy to technical enforcement through defined roles, procedures, and controls
\cite{elbaghazaoui_data_2025, elbaghazaoui_ensuring_2025}. Failures in either
dimension expose organisations to legal, financial, and reputational risk
\cite{elbaghazaoui_data_2025}.

DevOps emerged to accelerate software delivery through automation, continuous
integration and deployment (CI/CD), and rapid feedback cycles
\cite{elbaghazaoui_integrating_2025-1}. DevSecOps extends this model by embedding
security and compliance across the software development lifecycle. Practices
such as automated testing, compliance-as-code, artifact signing, and immutable
logging produce verifiable artifacts that can be inspected by regulators
\cite{anisetti_devsecops-based_2022, castellanos_accordant_2021,
elbaghazaoui_integrating_2025-1}. This convergence bridges engineering and
governance concerns, establishing compliance as a built-in property of the
delivery process.

As data itself becomes a primary strategic asset, surpassing the value
of the systems that process it \cite{elbaghazaoui_data_2025,
elbaghazaoui_ensuring_2025}, the mechanisms of DevSecOps must be extended to the
data pipelines that sustain organisational decision-making.

\section{Defining Key Concepts}
\noindent
The following subsections define the key concepts underpinning this research:
DevSecOps, ETL, data governance, and auditability. Each definition draws from
existing literature and sets the theoretical foundation for later analysis.

\subsection{DevSecOps}
\noindent
DevSecOps is the proactive and continuous enforcement of security and compliance
requirements throughout the software development lifecycle (SDLC)
\cite{anisetti_devsecops-based_2022}. Unlike traditional models that treat
security as an external stage or afterthought, DevSecOps embeds it within every
phase of the lifecycle \cite{anisetti_semi-automatic_2020}. Compliance becomes
the default state, and deviations from defined standards must be explicitly
justified \cite{castellanos_accordant_2021}.

A defining characteristic of DevSecOps is its reliance on automation to enforce
non-functional requirements such as auditability and governance. Through
automated testing, provisioning, monitoring, and policy enforcement, systems
produce immutable evidence—signed artifacts, logs, and metrics—that regulators
and auditors can independently assess \cite{zhou_towards_2020,
anisetti_devsecops-based_2022}. This automated assurance bridges the concerns of
development teams and governance stakeholders
\cite{anisetti_semi-automatic_2020}.

By combining cultural expectations of shared responsibility with technical
mechanisms for verification, DevSecOps establishes a delivery model in which
security and compliance are continuous, measurable, and intrinsic—rather than
externally imposed.

\subsection{ETL}
\noindent
Extract–Transform–Load (ETL) is a structured process that ingests, transforms,
and persists data into target repositories such as data warehouses, lakes, or
file stores \cite{kimball_data_2004}. It traditionally comprises three stages:
extraction, where data is sourced from heterogeneous systems; transformation,
where it is cleansed, standardised, integrated, and reshaped into coherent
semantic structures; and loading, where transformed data is stored for
downstream analysis \cite{kimball_data_2004}. Kimball and Caserta describe ETL
as both the most resource-intensive and most critical link between operational
data and analytical use.

Modern practice extends beyond this classical batch model. Approaches such as
extract–load–transform (ELT) and hybrid streaming–batch architectures leverage
distributed and cloud-native infrastructure \cite{mahmud_role_2022}. ETL now
encompasses not only data integration and quality but also governance and
compliance. In regulated industries, pipelines must guarantee lineage,
reproducibility, and auditability in addition to performance and scalability
\cite{mahmud_role_2022}.

ETL is thus best understood as a flexible socio-technical process. It concludes
when data is made persistent and reusable but intersects with broader lifecycle
concerns such as orchestration, monitoring, and policy enforcement. In this
sense, ETL pipelines underpin both analytical scalability and institutional
trust \cite{mahmud_role_2022}.

\subsection{Data Governance}
\noindent
Data governance provides the strategic framework that defines how data is
collected, managed, and protected across an organisation. It establishes the
policies, roles, standards, and processes that ensure information assets remain
accurate, consistent, secure, and compliant with internal and external
regulations \cite{elbaghazaoui_data_2025, elbaghazaoui_ensuring_2025}. Nai,
Rifai, and Sadiq~\cite{elbaghazaoui_data_2025} describe governance as a
structured framework that ensures the availability, integrity, usability, and
security of data through stewardship and accountability roles.

Vayyala~\cite{elbaghazaoui_ensuring_2025} expands this understanding by
positioning governance as the foundation of reliable data ecosystems—linking
policy to technical enforcement through automation, lineage tracking, and
immutable audit logs. In practice, governance institutionalises compliance and
accountability, transforming data into a verifiable strategic asset that can
withstand regulatory scrutiny.

\subsection{Auditability}
\noindent
Auditability is the designed capability of a system to allow its processes,
controls, and outcomes to be observed and verified by stakeholders. Weigand and
Elsas~\cite{weigand_auditability_2019} define auditability through five
interdependent requirements:
\begin{itemize}
  \item \textbf{Normativity}: a clear regulatory or contractual framework that
  governs operations.
  \item \textbf{Accountability}: verifiable statements linking actions to
  responsible agents and governing norms.
  \item \textbf{Referentiality Infrastructure}: records that connect operations
  to their original points of recording.
  \item \textbf{Reliability}: faithful and tamper-resistant representation of
  events.
  \item \textbf{Verifiability}: the ability of external parties to independently
  test or validate claims.
\end{itemize}

Although these principles originate in organisational and financial control,
they can be transferred to information systems. Extending this theoretical
foundation, Wang et al.~\cite{wang_enabling_2011} demonstrate how auditability
can be operationalised through third-party verification, cryptographic proofs of
integrity, and support for dynamic data operations. Their work illustrates that
auditability is not merely an organisational goal but a property that can be
engineered into data management infrastructures.

Applied to ETL pipelines and DevSecOps practices, these five dimensions map
directly to regulated-industry concerns:
\begin{itemize}
  \item \textbf{Normativity}: embedding compliance and data-protection policies
  within workflow execution.
  \item \textbf{Accountability}: linking each transformation or deployment to
  authenticated agents or automated processes.
  \item \textbf{Referentiality Infrastructure}: ensuring full lineage tracking
  across extraction, transformation, and loading.
  \item \textbf{Reliability}: maintaining immutable logs, signed artifacts, and
  cryptographic hashes.
  \item \textbf{Verifiability}: enabling reproducibility tests and independent
  audits via CI/CD-integrated checks.
\end{itemize}

Together, these adapted requirements frame auditability as a designed and
engineered property of ETL systems. Every data operation and system change must
be traceable, reproducible, and independently verifiable against regulatory
expectations.


\section{RQ1: ETL Pipelines in Regulated Industries}
\noindent

ETL pipelines in regulated industries operate under a dual mandate: they must
deliver analytical value while producing verifiable evidence of control. Mahmud
and Ikbal~\cite{mahmud_role_2022} describe ETL systems as socio-technical
infrastructures that extend beyond data integration to support institutional
trust. Their study highlights that scalability and performance alone are
insufficient; true maturity lies in sustaining lineage, transparency, and
compliance across distributed environments.

This tension is most visible in financial, healthcare, and critical
infrastructure domains, where frameworks such as APRA’s CPS~230, CPS~234, and
CPS~510~\cite{cps230,cps234,cps510} formalise separation of duties and
controlled change promotion. Broader governance research echoes these
principles, emphasising role-based accountability, immutable logging, and
retention aligned with compliance lifecycles
\cite{elbaghazaoui_ensuring_2025, elbaghazaoui_metadata_2025,
elbaghazaoui_data_2025}. Cloud-native orchestration services such as AWS Glue
and Azure Data Factory address elasticity demands, yet their adoption is
tempered by the need for traceability, access control, and evidence retention
consistent with CPG~235
\cite{australian_prudential_regulation_authority_prudential_2013}. The result
is a tightly controlled ecosystem where orchestrators such as Airflow and
serverless triggers coordinate data movement while preserving verifiable
execution histories~\cite{pogiatzis_event-driven_2021}.

Within this environment, ETL pipelines increasingly act as vehicles for
automated governance. Vayyala~\cite{elbaghazaoui_ensuring_2025} argues that
integrating DevSecOps-style validation—continuous testing, versioning, and
policy enforcement—transforms pipelines into compliance mechanisms rather than
mere conduits for data. Metadata management reinforces this shift. By capturing
provenance, transformation logic, and operational context, metadata forms the
connective tissue between policy and process
\cite{elbaghazaoui_metadata_2025}. Process metadata in particular documents each
transformation step, producing an intrinsic audit trail that enables
reproducibility and accountability.

Despite these advances, persistent obstacles limit complete auditability.
Heterogeneous toolchains fragment lineage across orchestration and storage
layers, while transformation logic embedded in SQL or user-defined functions
remains difficult to trace~\cite{mahmud_role_2022}. Schema drift and
inconsistent metadata propagation further erode reliability, and few systems
produce immutable, cryptographically verifiable artifacts by default
\cite{elbaghazaoui_metadata_2025, weigand_auditability_2019}. As Weigand and
Elsas~\cite{weigand_auditability_2019} note, such weaknesses disrupt the chain
of evidence essential to verifiable governance.

In summary, ETL workflows in regulated industries have evolved from operational
utilities into compliance-critical systems, embedding automation, lineage, and
policy enforcement at their core. Yet fragmented lineage, limited immutability,
and opaque transformation logic continue to hinder full auditability. These
gaps highlight the need for lifecycle-aware approaches that integrate security
and verification from design through deployment.

The next section examines how
DevSecOps practices—originating within the software development lifecycle—extend
these governance capabilities into data pipelines, bridging software assurance
and data lifecycle management.



\section{RQ2: DevSecOps and Governance in the SDLC}
\noindent

The software development lifecycle (SDLC) provides a disciplined framework for
the design, construction, and maintenance of software systems. It translates
functional and non-functional requirements into verifiable artefacts and
operational releases, maintaining quality, reliability, and security across
iterations. International standards such as ISO/IEC/IEEE~12207 and~15288
formalise these stages through defined information items, plans, and reports
that collectively establish traceability and accountability within engineering
processes \cite{noauthor_isoiecieee_2019, hernandez_requirements_2023}. The
SDLC thereby acts as a governance mechanism—an architecture of control in which
security, assurance, and auditability can be systematically embedded
\cite{souppaya_secure_2022, anisetti_devsecops-based_2022}.

As established in Chapter~2, auditability extends beyond retrospective review to
a system’s designed capacity to produce immutable, contextualised evidence of
operation. The SDLC becomes the medium through which this evidence is
constructed and preserved, transforming the development process itself into a
chain of accountability.

DevSecOps elevates this framework by embedding compliance and verification
directly within the lifecycle. Instead of treating assurance as post-hoc
oversight, DevSecOps positions it as a continuous, measurable property of
delivery. Machine-readable artefacts—signed binaries, immutable logs, and
compliance manifests—are generated automatically at each stage of development
\cite{anisetti_devsecops-based_2022, anisetti_semi-automatic_2020,
chen_design_2022, weigand_auditability_2019}. These artefacts form a living
record of integrity that binds design intent to operational reality.

In this model, pipelines are annotated with non-functional requirements that are
continuously verified through automated checks. Each stage yields testable
outputs that contribute to an operational chain of trust, supported by
continuous certification frameworks validating service behaviour against
assurance policies \cite{anisetti_semi-automatic_2020, chen_design_2022,
souppaya_secure_2022}. Auditability thus shifts from a retrospective activity to
a real-time system property, maintained by automation and policy-as-code.

Infrastructure-as-code (IaC), continuous integration and delivery (CI/CD), and
policy-as-code collectively translate this philosophy into practice. Static and
dynamic analysis, dependency scanning, and pre-deployment gates transform
traditional control functions into machine-verifiable checkpoints distributed
throughout the pipeline \cite{li_evolution_2025,
anisetti_semi-automatic_2020}. Governance artefacts emerge alongside the code
itself, enhancing reproducibility and traceability across environments
\cite{anisetti_devsecops-based_2022, hernandez_requirements_2023}. This
shift-left paradigm ensures that compliance is achieved through construction,
not inspection.

These practices mirror the Secure Software Development Framework (SSDF)
outlined in NIST Special Publication~800-218 \cite{souppaya_secure_2022}, which
defines four domains: Prepare the Organization, Protect Software, Produce
Well-Secured Software, and Respond to Vulnerabilities. NIST prescribes
traceable metadata, provenance records, and policy enforcement as essential
lifecycle artefacts—principles that align closely with DevSecOps’ focus on
continuous verification and toolchain provenance
\cite{anisetti_devsecops-based_2022, chen_design_2022}. When embedded within
CI/CD pipelines, these practices create persistent audit readiness and
measurable compliance, satisfying both technical and regulatory governance
requirements \cite{australian_prudential_regulation_authority_prudential_2013,
weigand_auditability_2019}.

Contemporary security research reinforces this need for governance-by-design.
The evolution of the OWASP Top~10 and CWE/SANS Top~25 demonstrates the rise of
supply-chain, API, cloud, and AI/ML vulnerabilities
\cite{li_evolution_2025}. Sustaining auditability under these conditions
requires continuous dependency scanning, provenance validation, and policy-as-
code \cite{souppaya_secure_2022, anisetti_devsecops-based_2022}. In regulated
domains such as automotive software, integrated security audits at code, build,
and runtime achieve real-time compliance validation while maintaining delivery
velocity \cite{achuthan_shifting_2024, chen_design_2022}. These architectures
prove that auditability and agility need not be in opposition.

Emerging techniques extend this assurance through intelligent automation.
Large language models (LLMs) are being embedded into enterprise pipelines to
interpret and enforce compliance policies. Hybrid systems combining rule-based
and AI-driven validation reduce manual review effort while maintaining
regulatory precision, enabling adaptive monitoring across Kubernetes,
Terraform, and IAM configurations \cite{mittal_practical_2025,
souppaya_secure_2022, anisetti_semi-automatic_2020}. Such augmentation
illustrates DevSecOps’ ability to evolve as a semantic layer of governance that
adapts with the technologies it protects.

Traceability remains the connective tissue binding these mechanisms together.
Integrated requirements management and automated linkage validation establish a
bidirectional audit trail between regulatory intent and deployed artefacts
\cite{noauthor_isoiecieee_2019, anisetti_devsecops-based_2022,
souppaya_secure_2022}. This alignment bridges organisational policy and
technical execution, ensuring that governance is both demonstrable and durable
\cite{hernandez_requirements_2023, chen_design_2022}. In cloud-native contexts,
where distributed workloads complicate oversight, embedded vulnerability
scanning, configuration policy enforcement, and runtime monitoring preserve
provenance and certification integrity
\cite{elbaghazaoui_integrating_2025-1, anisetti_devsecops-based_2022,
souppaya_secure_2022}. Every deployment becomes reconstructable, verifiable, and
provable in context.

These lifecycle-integrated practices parallel the expectations of APRA’s
Prudential Practice Guide CPG~235: Managing Data Risk
\cite{australian_prudential_regulation_authority_prudential_2013}, which
requires continuous validation, classification, and assurance of information
assets throughout their operational life. DevSecOps provides the technical means
to satisfy these expectations through automated control validation and evidence
generation \cite{anisetti_semi-automatic_2020,
weigand_auditability_2019}.  

At its core, auditability is a design property. Systems must be architected for
observability, referential integrity, and completeness of evidence
\cite{weigand_auditability_2019}. In a DevSecOps context, the CI/CD pipeline
itself functions as the referential infrastructure linking policy to artefact:
a dynamic chain of evidence that binds governance to execution
\cite{anisetti_devsecops-based_2022, souppaya_secure_2022,
chen_design_2022}. Auditability thus emerges not as a final checklist but as an
inherent characteristic of the lifecycle architecture
\cite{anisetti_semi-automatic_2020}.

Collectively, these developments show that DevSecOps transforms governance from
a procedural burden into a continuous property of software production. By
combining assurance, policy-as-code, and machine-verifiable artefacts, it fuses
engineering and compliance into a unified system of accountability—capable of
sustaining trust and regulatory transparency in real time
\cite{souppaya_secure_2022, anisetti_devsecops-based_2022,
anisetti_semi-automatic_2020, weigand_auditability_2019,
australian_prudential_regulation_authority_prudential_2013}.

The next section compares the data and software lifecycles, examining how this
integration of assurance can extend from code to data, bridging lifecycle-based
auditability into the domain of data governance.



\section{RQ3: Comparing the Data Lifecycle and the SDLC}
\noindent


\section{RQ4: Tooling for Auditability and Governance}
% Address sub-question 4
% - Categories of DevSecOps tools: CI/CD automation, IaC security, artifact signing, monitoring
% - Categories of ETL tools: orchestration (Airflow), processing (Spark), storage (S3), metadata management, lineage
% - Comparative discussion of features relevant to auditability and governance
% - Gaps where ETL tooling lacks equivalent DevSecOps maturity

\section{Synthesis}
% Summarize insights from RQ1–RQ4
% Highlight where existing work falls short of integrating DevSecOps with ETL
% Establish the research gap this thesis addresses
% Transition to the methodology chapter